\input{../output/tables/q2_numbers.tex}
\begin{abstract}
针对含光伏与储能的微网购电调控问题，本文以负荷满足为前提，将储能跨时段配置、预测不足风险和购电计划调整纳入统一模型。以储能电量和有效购电合同描述系统状态，以决策时可用信息限制策略，分别建立确定性日前计划、不确定条件下的固定合同、预报更新下的滚动调整及实时价格扩展模型。

\textbf{针对问题一，}建立包含能量平衡、储能动态、充放电互斥及日末电量恢复约束的混合整数线性规划，采用费用优先的分层求解。全天购电量为 $59482.699\ \mathrm{kWh}$，购电费为 $35126.85$ 元，相较无储能节费 $26.90\%$，且无弃光。整数最优值与线性松弛下界一致，最终方案通过独立物理及费用核验。

\textbf{针对问题二，}采用历史趋势预测与样本外残差分位数描述净负荷风险，仅根据一月费用选择日前规划及储能控制参数。最终采用80\%风险分位数和贪心储能缓冲，二月至十二月实际总费用为\QTwoTotalCost 万元，较同信息无储能基线节费\QTwoSaving\%。固定合同下逐时响应，跨日SOC连续，尾部风险另行报告。

\textbf{针对问题三，}在四个预报发布时点更新未来信息与实际储能状态，对尚未交付的合同进行滚动优化。分别核算原计划、上调、下调及紧急购电费用，以包含调整代价的净信息价值评价预报更新，并讨论不同下调结算解释的影响。

\textbf{针对问题四，}将未来价格纳入负荷与光伏的联合条件分布，分别扩展固定合同和可调整合同两类策略，保持设备约束与交易权限一致，区分用于规划的价格预测和用于结算的实际价格。

【待补充：问题三、四的实现、数值检验及新增预报时点结论。】

\keywords{微网调度；储能配置；非对称成本；滚动优化；信息约束}
\end{abstract}
